\documentclass[10pt,a4paper]{amsart}
\usepackage{amsmath,amssymb,amsthm,graphicx}
\usepackage{verbatim}
\newenvironment{reference}{\comment}{\endcomment}
%\newenvironment{reference}{}{}
\newenvironment{slogan}{\comment}{\endcomment}
\newenvironment{history}{\comment}{\endcomment}
% For commutative diagrams we use Xy-pic
\usepackage[all]{xy}

% We use 2cell for 2-commutative diagrams.
\xyoption{2cell}
% Theorem environments.
%
\theoremstyle{plain}
\newtheorem{theorem}[subsection]{Theorem}
\newtheorem{proposition}[subsection]{Proposition}
\newtheorem{lemma}[subsection]{Lemma}

\theoremstyle{definition}
\newtheorem{definition}[subsection]{Definition}
\newtheorem{example}[subsection]{Example}
\newtheorem{exercise}[subsection]{Exercise}
\newtheorem{situation}[subsection]{Situation}

\theoremstyle{remark}
\newtheorem{remark}[subsection]{Remark}
\newtheorem{remarks}[subsection]{Remarks}

\numberwithin{equation}{subsection}

% Macros
%
\def\lim{\mathop{\mathrm{lim}}\nolimits}
\def\colim{\mathop{\mathrm{colim}}\nolimits}
\def\Spec{\mathop{\mathrm{Spec}}}
\def\Hom{\mathop{\mathrm{Hom}}\nolimits}
\def\Ext{\mathop{\mathrm{Ext}}\nolimits}
\def\SheafHom{\mathop{\mathcal{H}\!\mathit{om}}\nolimits}
\def\SheafExt{\mathop{\mathcal{E}\!\mathit{xt}}\nolimits}
\def\Sch{\mathit{Sch}}
\def\Mor{\mathop{\mathrm{Mor}}\nolimits}
\def\Ob{\mathop{\mathrm{Ob}}\nolimits}
\def\Sh{\mathop{\mathit{Sh}}\nolimits}
\def\NL{\mathop{N\!L}\nolimits}
\def\CH{\mathop{\mathrm{CH}}\nolimits}
\def\proetale{{pro\text{-}\acute{e}tale}}
\def\etale{{\acute{e}tale}}
\def\QCoh{\mathit{QCoh}}
\def\Ker{\mathop{\mathrm{Ker}}}
\def\Im{\mathop{\mathrm{Im}}}
\def\Coker{\mathop{\mathrm{Coker}}}
\def\Coim{\mathop{\mathrm{Coim}}}

% Boxtimes
%
\DeclareMathSymbol{\boxtimes}{\mathbin}{AMSa}{"02}

%
% Macros for moduli stacks/spaces
%
\def\QCohstack{\mathcal{QC}\!\mathit{oh}}
\def\Cohstack{\mathcal{C}\!\mathit{oh}}
\def\Spacesstack{\mathcal{S}\!\mathit{paces}}
\def\Quotfunctor{\mathrm{Quot}}
\def\Hilbfunctor{\mathrm{Hilb}}
\def\Curvesstack{\mathcal{C}\!\mathit{urves}}
\def\Polarizedstack{\mathcal{P}\!\mathit{olarized}}
\def\Complexesstack{\mathcal{C}\!\mathit{omplexes}}
% \Pic is the operator that assigns to X its picard group, usage \Pic(X)
% \Picardstack_{X/B} denotes the Picard stack of X over B
% \Picardfunctor_{X/B} denotes the Picard functor of X over B
\def\Pic{\mathop{\mathrm{Pic}}\nolimits}
\def\Picardstack{\mathcal{P}\!\mathit{ic}}
\def\Picardfunctor{\mathrm{Pic}}
\def\Deformationcategory{\mathcal{D}\!\mathit{ef}}
\usepackage{cleveref}
\crefname{originalreference}{original reference}{original references}
\Crefname{originalreference}{Original reference}{Original references}
\makeatletter
\newlabel{algebra-lemma-intersection-powers-ideal-module}{{algebra-lemma-intersection-powers-ideal-module}{source}{}{sourcepointer-1}{}}
\newlabel{algebra-lemma-intersection-powers-ideal-module@cref}{{[originalreference][0][]algebra-lemma-intersection-powers-ideal-module}{[1][0][]source}{}{}{}}
\makeatother
\crefname{originallemma}{original lemma}{original lemmas}
\Crefname{originallemma}{Original lemma}{Original lemmas}
\makeatletter
\newlabel{lemma-completion-local}{{lemma-completion-local}{source}{}{sourcepointer-2}{}}
\newlabel{lemma-completion-local@cref}{{[originallemma][0][]lemma-completion-local}{[1][0][]source}{}{}{}}
\makeatother
\title{Description of the zeroth local cohomology of the completion as a kernel of localizations}
\author{The Stacks Project Authors}
\date{stacks-project:dualizing.tex:local-snapshot-20261012; GNU FDL 1.2 or later}
\begin{document}
\maketitle
\noindent\textit{Editorial note.} All delivered mathematical text is the author's own, copied from the pinned Stacks Project chapter source sources/stacks-dualizing/pinned-original-source.tex; no statement or proof step was rewritten, reordered or removed.\par\medskip
\noindent\textit{Editorial note.} Preamble provenance: the Stacks Project keeps one single shared preamble for every chapter; this card takes it directly from the staged sources/stacks-dualizing/preamble.tex as three verbatim spans (lines 8-12, 14-18 and 154-217), so no derived file was written and nothing was added or redefined. The wrapper class and the base packages are supplied by the assembler.\par\medskip
\noindent\textit{Editorial note.} Bibliography provenance: neither the counted statement nor its proof contains a citation (this lane rescans the whole span before writing the card), so the card carries no bibliography and no entry was reconstructed.\par\medskip
\noindent\textit{Editorial note.} Selection provenance: this card was opened under the file-based control contract (CONTROL.md), with the mechanical pre-card gate run on its statement and proof interval before the card was written.\par\medskip
\noindent\textit{Allowed standard prerequisites (editorial).}\begin{itemize}
\item The description of the zeroth local cohomology of the completion as a kernel of localised completions from the proof of the local completion lemma. Invoked at the start of the counted proof (source lines 2381-2387). Supplies the displayed identification.
\item The statement on the vanishing of the intersection of the powers of an ideal after localising. Invoked inside the counted proof (source lines 2390-2393) as a cross-chapter pointer. Shows the local vanishing of the kernel pieces.
\end{itemize}\medskip
\ifdefined\hypertarget\hypertarget{sourcepointer-1}{}\fi
\noindent\textit{Source reference.} Original reference algebra-lemma-intersection-powers-ideal-module: sources/stacks-dualizing/pinned-original-source.tex, source lines 2390--2390. Navigation pointer only: the counted text refers to the Stacks item algebra-lemma-intersection-powers-ideal-module, which belongs to another chapter of the project and is not part of this card. The pointer carries the target's own label as its number because this chapter's staged snapshot does not contain the other chapter's tag table; its content is not reproduced here.\par\medskip
\ifdefined\hypertarget\hypertarget{sourcepointer-2}{}\fi
\noindent\textit{Source reference.} Original lemma lemma-completion-local: sources/stacks-dualizing/pinned-original-source.tex, source lines 2327--2327. Navigation pointer only: the counted text refers to the Stacks item lemma-completion-local, which is an existing labeled result of this same chapter and is not part of this card. The pointer carries the target's own label as its number and its definition line as its span; its content is not reproduced here.\par\medskip
\begin{lemma}
\label{lemma-completion-local-H0}
With notation and hypotheses as in Lemma \ref{lemma-completion-local}
assume $A$ is $I$-adically complete. Then
$$
H^0(R\Gamma_Z(M)^\wedge) = \colim H^0_{V(J')}(M)
$$
where the filtered colimit is over $J' \subset J$ such that
$V(J') \cap V(I) = V(J) \cap V(I)$.
\end{lemma}

\begin{proof}
Since $M$ is a finite $A$-module, we have that $M$ is $I$-adically complete.
The proof of Lemma \ref{lemma-completion-local} shows that
$$
H^0(R\Gamma_Z(M)^\wedge) =
\Ker(M^\wedge \to \prod M_{g_j}^\wedge) =
\Ker(M \to \prod M_{g_j}^\wedge)
$$
where on the right hand side we have usual $I$-adic completion.
The kernel $K_j$ of $M_{g_j} \to M_{g_j}^\wedge$ is $\bigcap I^n M_{g_j}$.
By Algebra, Lemma \ref{algebra-lemma-intersection-powers-ideal-module}
for every $\mathfrak p \in V(IA_{g_j})$ we find an
$f \in A_{g_j}$, $f \not \in \mathfrak p$ such that $(K_j)_f = 0$.

\medskip\noindent
Let $s \in H^0(R\Gamma_Z(M)^\wedge)$.
By the above we may think of $s$ as an element of $M$.
The support $Z'$ of $s$ intersected with $D(g_j)$ is disjoint from
$D(g_j) \cap V(I)$ by the arguments above.
Thus $Z'$ is a closed subset of $\Spec(A)$ with $Z' \cap V(I) \subset V(J)$.
Then $Z' \cup V(J) = V(J')$ for some ideal $J' \subset J$ with
$V(J') \cap V(I) \subset V(J)$ and we have $s \in H^0_{V(J')}(M)$.
Conversely, any $s \in H^0_{V(J')}(M)$ with $J' \subset J$ and
$V(J') \cap V(I) \subset V(J)$ maps to zero in $M_{g_j}^\wedge$ for all $j$.
This proves the lemma.
\end{proof}

\end{document}
